% =========================================================
% format_exam.tex
% Exam / test layout compatible with preamble_1.tex
% Style direction aligned with format_2.tex
% Target compiler: XeLaTeX
% =========================================================

% -------------------------
% 0) Extra packages not already in preamble_1
% -------------------------
\RequirePackage{xparse}
\RequirePackage{etoolbox}
\RequirePackage{calc}
\RequirePackage{lastpage}
\tcbuselibrary{skins,breakable}

% -------------------------
% 1) Colors / visual identity
% -------------------------
% Reuse existing palette names from format_2 when possible.
% If format_2 is not loaded, define them here.
\providecolor{GDBlue}{HTML}{127BBE}
\providecolor{GDDeepBlue}{HTML}{0000FF}
\providecolor{GDPink}{HTML}{C1007A}
\providecolor{GDOrange}{HTML}{C8741D}
\providecolor{GDPurple}{HTML}{6D2C91}
\providecolor{GDTeal}{HTML}{0E8C8C}
\providecolor{GDLiteBlue}{HTML}{EAF5FF}
\providecolor{GDLiteGray}{HTML}{F5F7FA}

% -------------------------
% 2) Local helper lengths (safe defaults if not declared)
% -------------------------
\providecommand{\GDExamHeaderRule}{1.0pt}
\newlength{\GDExamTitleSep}      \setlength{\GDExamTitleSep}{0.05cm}
\newlength{\GDExamBlockSep}      \setlength{\GDExamBlockSep}{0.05cm}
\newlength{\GDExamChoiceSep}     \setlength{\GDExamChoiceSep}{10pt}
\newlength{\GDExamChoiceRowSep}  \setlength{\GDExamChoiceRowSep}{5pt}
\newlength{\GDExamPartAfterSkip} \setlength{\GDExamPartAfterSkip}{0pt}
\newlength{\GDExamQuestionSkip}  \setlength{\GDExamQuestionSkip}{0pt}

% -------------------------
% 3) Small inline tag helper
% -------------------------
\providecommand{\GDTag}[2]{%
  \tikz[baseline=(n.base)]\node (n) [fill=#1, text=white, font=\bfseries\footnotesize,
    rounded corners=2pt, inner xsep=6pt, inner ysep=2pt] {#2};%
}

% -------------------------
% 4) Header information
% -------------------------
\makeatletter
\newcommand{\SchoolName}[1]{\gdef\GD@SchoolName{#1}}
\newcommand{\ExamTitle}[1]{\gdef\GD@ExamTitle{#1}}
\newcommand{\SubjectName}[1]{\gdef\GD@SubjectName{#1}}
\newcommand{\ExamDuration}[1]{\gdef\GD@ExamDuration{#1}}
\newcommand{\ExamNote}[1]{\gdef\GD@ExamNote{#1}}
\newcommand{\ExamCode}[1]{\gdef\GD@ExamCode{#1}}

\newcommand{\GD@SchoolName}{}
\newcommand{\GD@ExamTitle}{}
\newcommand{\GD@SubjectName}{}
\newcommand{\GD@ExamDuration}{}
\newcommand{\GD@ExamNote}{}
\newcommand{\GD@ExamCode}{}
\newcommand{\GD@StudentNameLabel}{Họ và tên}
\newcommand{\GD@StudentIDLabel}{Số báo danh}

% -------------------------
% Global mode: exam / solution
% -------------------------
\newif\ifshowanswers
\showanswersfalse % default = exam version

\newcommand{\ExamMode}{\showanswersfalse}
\newcommand{\SolutionMode}{\showanswerstrue}

% -------------------------
% Configurable labels and QR behavior
% -------------------------
\newcommand{\SetSolutionTitle}[1]{\gdef\GD@SolutionTitle{#1}}
\newcommand{\SetAnswerLabel}[1]{\gdef\GD@AnswerLabel{#1}}
\newcommand{\SetTrueSymbol}[1]{\gdef\GD@TrueSymbol{#1}}
\newcommand{\SetFalseSymbol}[1]{\gdef\GD@FalseSymbol{#1}}

\newcommand{\GD@SolutionTitle}{Hướng dẫn giải}
\newcommand{\GD@AnswerLabel}{Đáp số:}
\newcommand{\GD@TrueSymbol}{Đ}
\newcommand{\GD@FalseSymbol}{S}

\newif\ifGDShowHeaderQR
\GDShowHeaderQRfalse
\newcommand{\ShowHeaderQR}{\GDShowHeaderQRtrue}
\newcommand{\HideHeaderQR}{\GDShowHeaderQRfalse}

\newif\ifGDShowExamHeader
\GDShowExamHeadertrue
\newcommand{\ShowExamHeader}{\GDShowExamHeadertrue}
\newcommand{\HideExamHeader}{\GDShowExamHeaderfalse}

\newcommand{\ExamQR}[1]{\gdef\GD@ExamQR{#1}}
\newcommand{\GD@ExamQR}{}
\newcommand{\GDQRWidth}{2.6cm}

% A single command to set all fields quickly
\NewDocumentCommand{\ExamInfo}{m m m m m m}{%
  \SchoolName{#1}%
  \ExamTitle{#2}%
  \SubjectName{#3}%
  \ExamDuration{#4}%
  \ExamNote{#5}%
  \ExamCode{#6}%
}

% -------------------------
% 5) Exam title block
% -------------------------
\newcommand{\MakeExamHeader}{%
  \ifGDShowExamHeader
    \par\noindent\small
    \begin{minipage}[t]{0.24\linewidth}
      \vspace{0pt}
      \centering
      {\bfseries \GD@SchoolName\par}
    \end{minipage}
    \hfill
    \begin{minipage}[t]{0.56\linewidth}
      \vspace{0pt}
      \centering
      {\bfseries \GD@ExamTitle\par}
      \vspace{2pt}
      {\bfseries Môn: \GD@SubjectName\par}
      \vspace{2pt}
      Thời gian làm bài: \GD@ExamDuration\ifstrempty{\GD@ExamNote}{}{, \GD@ExamNote}.\par
      \vspace{2pt}
      \ifstrempty{\GD@ExamCode}{}{{\bfseries Mã đề: \GD@ExamCode\par}}
    \end{minipage}
    \hfill
    \begin{minipage}[t]{0.16\linewidth}
      \vspace{0pt}
      \centering
        \ifGDShowHeaderQR
          \ifstrempty{\GD@ExamQR}{}{%
            \includegraphics[width=\GDQRWidth]{\GD@ExamQR}\par
          }
        \fi
    \end{minipage}
    \par\vspace{\GDExamBlockSep}%
  \fi
}

% Optional student information line
\newcommand{\MakeStudentInfoLine}{%
  \par\noindent\normalsize \GD@StudentNameLabel: \dotfill \hfill \GD@StudentIDLabel: \dotfill \par
  \vspace{\GDExamBlockSep}%
}
\makeatother

% -------------------------
% 6) Counters
% -------------------------
\newcounter{examPart}
\newcounter{examQuestion}[examPart]

\renewcommand{\theexamPart}{\arabic{examPart}}
\renewcommand{\theexamQuestion}{\arabic{examQuestion}}

% utility setters
\newcommand{\SetExamPartStart}[1]{\setcounter{examPart}{\numexpr#1-1\relax}}
\newcommand{\SetExamQuestionStart}[1]{\setcounter{examQuestion}{\numexpr#1-1\relax}}

% -------------------------
% 7) Part heading
% Requirement: auto numbering + configurable start + points + description
% Usage:
%   \ExamPart{Phần trắc nghiệm}{3,0}{Chọn đáp án đúng nhất.}
%   \ExamPart[5]{Phần tự luận}{7,0}{Trình bày đầy đủ lời giải.}
% -------------------------
\NewDocumentCommand{\ExamPart}{O{} m m m}{%
  \par\vspace{\GDExamBlockSep}%
  \ifstrempty{#1}{}{%
    \setcounter{examPart}{\numexpr#1-1\relax}%
  }%
  \refstepcounter{examPart}%
  \setcounter{examQuestion}{0}%
  \noindent\textbf{Phần \theexamPart.\ \textit{#2}\ (#3 điểm)}\par
  \noindent\textit{#4}\par
  \vspace{\GDExamPartAfterSkip}%
}

% -------------------------
% 8) Question heading
% Requirement: auto numbering
% Usage:
%   \Question Nội dung...
% -------------------------
\NewDocumentCommand{\Question}{}{%
  \par\vspace{\GDExamQuestionSkip}%
  \refstepcounter{examQuestion}%
  \noindent{\bfseries\color{GDDeepBlue}Câu \theexamQuestion. }\ignorespaces
}

% Robust alternative to \Question (uses \textbf + \textcolor instead of \bfseries\color)
\NewDocumentCommand{\NQuestion}{}{%
  \par\vspace{\GDExamQuestionSkip}%
  \refstepcounter{examQuestion}%
  \noindent\textbf{\textcolor{GDDeepBlue}{Câu~\theexamQuestion.}}\ \ignorespaces
}

% Optional question line with score tag if needed later
\NewDocumentCommand{\QuestionWithPoints}{m}{%
  \par\vspace{\GDExamQuestionSkip}%
  \refstepcounter{examQuestion}%
  \noindent{\bfseries\color{GDDeepBlue}Câu \theexamQuestion. }\ignorespaces
  \hfill {\GDTag{GDOrange!85!black}{#1 điểm}}\par
  \noindent\ignorespaces
}

% -------------------------
% 9) a), b), c), d) sub-questions
% Requirement: auto letter increase + configurable start (default a)
% Usage:
%   \begin{SubQuestions}
%     \item ...
%   \end{SubQuestions}
%
%   \begin{SubQuestions}[3] % starts at c)
%     \item ...
%   \end{SubQuestions}
% -------------------------
\NewDocumentEnvironment{SubQuestions}{O{1}}{%
  \begin{enumerate}[
    label=\alph*),
    leftmargin=2.1em,
    itemsep=\GDListItemSep,
    topsep=\GDListTopSep
  ]
  \setcounter{enumi}{\numexpr#1-1\relax}
}{%
  \end{enumerate}
}

% -------------------------
% 10) Multiple choice labels A. B. C. D.
% Requirement: 3 layouts
% - 4 choices on 1 line
% - 2 choices on each line, 2 lines total
% - 1 choice per line, 4 lines total
% -------------------------
\newcounter{GDExamChoiceCounter}
\newcommand{\ResetExamChoices}{\setcounter{GDExamChoiceCounter}{0}}
\newcommand{\ExamChoiceLabel}{\stepcounter{GDExamChoiceCounter}\textbf{\Alph{GDExamChoiceCounter}.}}
\newcommand{\ExamChoiceText}[1]{\ExamChoiceLabel\,#1}

% -------------------------
% Answer highlight and solution helpers
% -------------------------
\providecolor{GDAnswerGreen}{HTML}{D9F2D9}
\providecolor{GDAnswerBorder}{HTML}{4CAF50}

\newcommand{\AnswerHighlight}[1]{%
  \tikz[baseline=(n.base)]\node[
    fill=GDAnswerGreen,
    draw=GDAnswerBorder,
    rounded corners=2pt,
    inner xsep=3pt,
    inner ysep=2pt
  ] (n) {#1};%
}

\newcommand{\GDIfEqIgnoreCase}[4]{%
  \begingroup
  \edef\GDTempA{\expandafter\lowercase\expandafter{#1}}%
  \edef\GDTempB{\expandafter\lowercase\expandafter{#2}}%
  \ifdefstrequal{\GDTempA}{\GDTempB}{\endgroup#3}{\endgroup#4}%
}

\newcommand{\GDChoiceLabelFromIndex}[1]{%
  \ifcase#1\relax ?\or A\or B\or C\or D\else ?\fi
}

\newcommand{\GDChoiceItem}[3]{%
  % #1 = choice index (1..4)
  % #2 = displayed text
  % #3 = optional correct index (1..4), empty = no highlight
  \ifshowanswers
    \ifstrempty{#3}{%
      \textbf{\GDChoiceLabelFromIndex{#1}.}~#2%
    }{%
      \ifnum#1=#3\relax
        \AnswerHighlight{\textbf{\GDChoiceLabelFromIndex{#1}.}~#2}%
      \else
        \textbf{\GDChoiceLabelFromIndex{#1}.}~#2%
      \fi
    }%
  \else
    \textbf{\GDChoiceLabelFromIndex{#1}.}~#2%
  \fi
}

\NewDocumentCommand{\SolutionBlock}{m}{%
  \ifshowanswers
    \par\smallskip
    \begin{center}
      \noindent\textbf{\csname GD@SolutionTitle\endcsname}\par
    \end{center}
    \noindent #1\par
  \fi
}

\NewDocumentCommand{\AnswerLine}{m}{%
  \ifshowanswers
    \par\smallskip
    \noindent\textbf{\csname GD@AnswerLabel\endcsname\ }\AnswerHighlight{#1}\par
  \fi
}

\newcommand{\GDTrueFalseSymbol}[1]{%
  \GDIfEqIgnoreCase{#1}{true}{\csname GD@TrueSymbol\endcsname}{\csname GD@FalseSymbol\endcsname}%
}

\NewDocumentCommand{\TFItem}{m m}{%
  % #1 = true/false (case-insensitive)
  % #2 = statement text
  \item #2%
  \ifshowanswers
    \hfill \AnswerHighlight{\textbf{\GDTrueFalseSymbol{#1}}}%
  \fi
}

% 4 choices on one line, optional answer letter as 5th braced arg
\newcommand{\GDChoicesOneLineWithAnswer}[5]{%
  \par\ResetExamChoices
  \noindent
  \begin{tabularx}{\linewidth}{@{}>{\raggedright\arraybackslash}X@{\hspace{\GDExamChoiceSep}}>{\raggedright\arraybackslash}X@{\hspace{\GDExamChoiceSep}}>{\raggedright\arraybackslash}X@{\hspace{\GDExamChoiceSep}}>{\raggedright\arraybackslash}X@{}}
    \GDChoiceItem{1}{#1}{#5} & \GDChoiceItem{2}{#2}{#5} & \GDChoiceItem{3}{#3}{#5} & \GDChoiceItem{4}{#4}{#5}
  \end{tabularx}
  \par
}

\newcommand{\GDChoicesOneLineNoAnswer}[4]{%
  \par\ResetExamChoices
  \noindent
  \begin{tabularx}{\linewidth}{@{}>{\raggedright\arraybackslash}X@{\hspace{\GDExamChoiceSep}}>{\raggedright\arraybackslash}X@{\hspace{\GDExamChoiceSep}}>{\raggedright\arraybackslash}X@{\hspace{\GDExamChoiceSep}}>{\raggedright\arraybackslash}X@{}}
    \GDChoiceItem{1}{#1}{} & \GDChoiceItem{2}{#2}{} & \GDChoiceItem{3}{#3}{} & \GDChoiceItem{4}{#4}{}
  \end{tabularx}
  \par
}

\makeatletter
\newcommand{\ChoicesOneLine}[4]{%
  \@ifnextchar\bgroup
    {\GDChoicesOneLineWithAnswer{#1}{#2}{#3}{#4}}%
    {\GDChoicesOneLineNoAnswer{#1}{#2}{#3}{#4}}%
}
\makeatother

% 2 choices per line, 2 lines, optional answer letter as 5th braced arg
\newcommand{\GDChoicesTwoLinesWithAnswer}[5]{%
  \par\ResetExamChoices
  \noindent
  \begin{tabularx}{\linewidth}{@{}>{\raggedright\arraybackslash}X@{\hspace{\GDExamChoiceSep}}>{\raggedright\arraybackslash}X@{}}
    \GDChoiceItem{1}{#1}{#5} & \GDChoiceItem{2}{#2}{#5} \\[\GDExamChoiceRowSep]
    \GDChoiceItem{3}{#3}{#5} & \GDChoiceItem{4}{#4}{#5}
  \end{tabularx}
  \par
}

\newcommand{\GDChoicesTwoLinesNoAnswer}[4]{%
  \par\ResetExamChoices
  \noindent
  \begin{tabularx}{\linewidth}{@{}>{\raggedright\arraybackslash}X@{\hspace{\GDExamChoiceSep}}>{\raggedright\arraybackslash}X@{}}
    \GDChoiceItem{1}{#1}{} & \GDChoiceItem{2}{#2}{} \\[\GDExamChoiceRowSep]
    \GDChoiceItem{3}{#3}{} & \GDChoiceItem{4}{#4}{}
  \end{tabularx}
  \par
}

\makeatletter
\newcommand{\ChoicesTwoLines}[4]{%
  \@ifnextchar\bgroup
    {\GDChoicesTwoLinesWithAnswer{#1}{#2}{#3}{#4}}%
    {\GDChoicesTwoLinesNoAnswer{#1}{#2}{#3}{#4}}%
}
\makeatother

% 1 choice per line, 4 lines, optional answer letter as 5th braced arg
\newcommand{\GDChoicesFourLinesWithAnswer}[5]{%
  \par\ResetExamChoices
  \noindent
  \begin{tabularx}{\linewidth}{@{}>{\raggedright\arraybackslash}X@{}}
    \GDChoiceItem{1}{#1}{#5} \\[\GDExamChoiceRowSep]
    \GDChoiceItem{2}{#2}{#5} \\[\GDExamChoiceRowSep]
    \GDChoiceItem{3}{#3}{#5} \\[\GDExamChoiceRowSep]
    \GDChoiceItem{4}{#4}{#5}
  \end{tabularx}
  \par
}

\newcommand{\GDChoicesFourLinesNoAnswer}[4]{%
  \par\ResetExamChoices
  \noindent
  \begin{tabularx}{\linewidth}{@{}>{\raggedright\arraybackslash}X@{}}
    \GDChoiceItem{1}{#1}{} \\[\GDExamChoiceRowSep]
    \GDChoiceItem{2}{#2}{} \\[\GDExamChoiceRowSep]
    \GDChoiceItem{3}{#3}{} \\[\GDExamChoiceRowSep]
    \GDChoiceItem{4}{#4}{}
  \end{tabularx}
  \par
}

\makeatletter
\newcommand{\ChoicesFourLines}[4]{%
  \@ifnextchar\bgroup
    {\GDChoicesFourLinesWithAnswer{#1}{#2}{#3}{#4}}%
    {\GDChoicesFourLinesNoAnswer{#1}{#2}{#3}{#4}}%
}
\makeatother

% Lowercase aliases for convenience
\let\choiceoneline\ChoicesOneLine
\let\choicetwolines\ChoicesTwoLines
\let\choicefourlines\ChoicesFourLines

% -------------------------
% 11) Convenience wrappers for common exam sections
% -------------------------
\newcommand{\PartTracNghiem}[2]{\ExamPart{Trắc nghiệm}{#1}{#2}}
\newcommand{\PartDungSai}[2]{\ExamPart{Đúng -- Sai}{#1}{#2}}
\newcommand{\PartTuLuan}[2]{\ExamPart{Tự luận}{#1}{#2}}

% -------------------------
% Two-column layout command
% Usage: \TwoColumnLayout{left content}{right content}
% Creates two independent minipages side by side
% -------------------------
\newcommand{\TwoColumnLayout}[2]{%
  \par
  \begin{minipage}[t]{0.48\textwidth}
    #1
  \end{minipage}
  \hfill
  \begin{minipage}[t]{0.48\textwidth}
    #2
  \end{minipage}
  \par
}

% -------------------------
% 12) Optional boxed answer area / writing lines
% -------------------------
\newcommand{\ExamDottedLines}[1]{%
  \par\noindent
  \foreach \i in {1,...,#1}{%
    \noindent\tikz{\draw[GDDeepBlue!60, dotted, line width=\DottedLineWidth] (0,0) -- (\linewidth,0);}\[\DottedLineGap]
  }%
}

% -------------------------
% 13) Mini demo header/footer for exam pages
% -------------------------
\newcommand{\SetExamHeader}[2]{%
  \pagestyle{fancy}
  \fancyhf{}
  \fancyhead[L]{\scriptsize\color{GDDeepBlue}#2\;\textbf{\thepage/\pageref{LastPage}}}
  \fancyhead[R]{\scriptsize\color{GDDeepBlue}#1}
  \fancyfoot[R]{\scriptsize\color{GDDeepBlue}#2\;\textbf{\thepage/\pageref{LastPage}}}
  % Last page style with --- HẾT --- marker
  \fancypagestyle{lastexampage}{%
    \fancyhf{}
    \fancyhead[L]{\scriptsize\color{GDDeepBlue}#2\;\textbf{\thepage/\pageref{LastPage}}}
    \fancyhead[R]{\scriptsize\color{GDDeepBlue}#1}
    \fancyfoot[R]{\scriptsize\color{GDDeepBlue}#2\;\textbf{\thepage/\pageref{LastPage}}}
  }%
  \AtEndDocument{\thispagestyle{lastexampage}}%
}

\newcommand{\DottedLines}[1]{%
  \par\noindent
  \foreach \i in {1,...,#1}{%
    \noindent\tikz{\draw[GDBlue!70, dotted, line width=\DottedLineWidth] (0,0) -- (\linewidth,0);}\\[\DottedLineGap]
  }%
}
% -------------------------
% 14) Example usage (comment only)
% -------------------------
% \ExamInfo{THPT ABC}{ĐỀ KIỂM TRA GIỮA KÌ I -- NĂM HỌC 2025--2026}{Toán 10}{90 phút}[Đề có 02 trang]
% \SetExamHeader{Đề kiểm tra Toán 10}{Trang}
% \MakeExamHeader
% \MakeStudentInfoLine
%
% \ExamPart{Trắc nghiệm}{3}{Chọn đáp án đúng nhất cho mỗi câu sau.}
% \Question Cho hàm số \(f\left(x\right)=x^2\). Giá trị \(f(2)\) là:
% \ChoicesOneLine{1}{2}{4}{8}
%
% \Question Tập nghiệm của phương trình \(x^2=1\) là:
% \ChoicesTwoLines{\(\{-1;1\}\)}{\(\{1\}\)}{\(\{-1\}\)}{\(\varnothing\)}
%
% \ExamPart{Đúng -- Sai}{2}{Xét từng ý sau và xác định đúng hay sai.}
% \Question Với hàm số \(y=x^2\):
% \begin{SubQuestions}
%   \item Đồ thị đi qua gốc tọa độ.
%   \item Đồ thị nhận trục \(Oy\) làm trục đối xứng.
%   \item Hàm số nghịch biến trên \(\mathbb{R}\).
%   \item \(f(-1)=1\).
% \end{SubQuestions}
%
% \ExamPart[5]{Tự luận}{5}{Trình bày đầy đủ lời giải.}
% \Question Giải phương trình \(x^2-5x+6=0\).
% \begin{SubQuestions}[2]
%   \item Phân tích đa thức thành nhân tử.
%   \item Suy ra nghiệm của phương trình.
% \end{SubQuestions}
%
% % Two-column layout example:
% % \TwoColumnLayout{
% %   \NQuestion Question in left column.
% %   \DottedLines{3}
% % }{
% %   \NQuestion Question in right column.
% %   \DottedLines{3}
% % }

